\begin{abstract}
  In the search for robust continual learning methods, model merging has emerged as a promising paradigm to consolidate task-specific capabilities without parameter interference. A prominent recent proposal, Sharpness-Aware Isotropic Merging (SAIM), claims that a dual-stage approach—combining a coordinate-restricted sharpness-aware optimizer (SA-BCD) with an SVD-based adaptive isotropic merging algorithm—is essential for avoiding representation collapse. In this work, we present a rigorous methodological deconstruction of SAIM to evaluate its constituent parts and isolate the true causal drivers of merging performance. By executing a multi-axial grid of 15 configurations (crossing 5 optimizers and 3 merging strategies) on Split CIFAR-100 with a Vision Transformer, we reveal important conceptual nuances and boundary conditions. First, we show that under standard sequential fine-tuning parity (the $\lambda = 0$ regime), SVD-based isotropic merging is redundant and actively distorts un-mixed parameters, degrading performance compared to naive weight averaging. Second, we find that optimizer-driven flatness is the foundational driver of merging success: simply training experts with standard, globally perturbed Sharpness-Aware Minimization (SAM) and merging them via naive Task Arithmetic yields a 9.87\% accuracy improvement over standard baselines. Finally, we evaluate an active weight-mixing regime ($\lambda = 0.2$) and discover that while SVD-based isotropic merging acts as a beneficial regularizer against active parameter interference (boosting SAM to 76.42\%), unconstrained global flatness (SAM) remains the critical foundation for success. Our findings offer a balanced, scholarly audit of multi-component merging pipelines and outline simpler, superior optimization-driven design choices.
\end{abstract}
